\subsection{The Discriminant} 
The quadratic formula ``normally'' gives 2 solutions, but it may give 1 or 0, depending on the equation.
\begin{Example}
Use the quadratic formula to solve each quadratic equation.
\begin{lpicvsplit}{3}{7}
\foreach \x/\c in {0/3,1/4,2/7} {
	\regtext{0.5*\value{fcols}+\x*\value{fcols}+\x}{-1}{$x^2+4x+\c=0$};
}
\end{lpicvsplit}
\end{Example}
\noi The number of solutions depends on the number inside the square root. This number has a name.
\begin{definition}{Discriminant}
The \term{discriminant} of a quadratic equation $ax^2+bx+c=0$ is the number
$$D=b^2-4ac$$
whose square root appears in the quadratic formula. The quadratic equation has:
\begin{itemize*}
\item Two real solutions if $D>0$
\item One real solution if $D=0$
\item No real solutions if $D<0$
\end{itemize*}
If $D$ is a perfect square, the solutions are rational. Otherwise, they are irrational.
\end{definition}
\begin{Example}
For each quadratic equation below, use the discriminant to determine how many real solutions the equation has.
\begin{lpicvsplit}{3}{5}
\foreach \x/\exp in {0/3x^2-11x-5,1/t^2+5t+25,2/x^2+6x+9}
	\regtext{0.5*\value{fcols}+\x*\value{fcols}+\x}{-1}{$\exp=0$};
\end{lpicvsplit}
\end{Example}